% From source to a running simulation (report r02, chapter "Code Generation" and its figure
% build-flow.tex). CMake configures; during "make", CMake runs the Python generator scripts
% (sweepgen sets jit = no_jit: generate only), then compiles the generated C++ together with
% the hand-written waLBerla core. The executable runs without Python. The time loop runs the kernel, the halo
% exchange and the wall pass (generated index-list loops, Target.GenericCPU: plain C++ in every build). Everything ran on the
% board. Violet = generated code. Each column is one step; the small arrows between the
% headers carry the results from step to step (no wires across columns).
%
% With \brmarkerstrue (slide "Where each change lives") numbered markers show what we changed:
% blue = part 1, teal = part 2 (10 to 12: the last three changes, see the open-work slide).
% The markers sit on the nodes a2, g0, g2, g3, i1, i2 and cc.
\begin{tikzpicture}[
    hdr/.style={font=\scriptsize\bfseries,text=white,align=center,inner sep=2pt,
      minimum height=8mm,anchor=north west},
    b/.style={box,font=\scriptsize,minimum height=5.5mm,inner sep=2pt,
      execute at begin node={\hyphenpenalty=10000\exhyphenpenalty=10000}},
    gen/.style={b,fill=psviolet!12,draw=psviolet!80!black},
    hdrarr/.style={fill=psgray!80},
    mk/.style={circle,text=white,font=\tiny\bfseries,inner sep=0pt,minimum size=3.7mm}]
  \ifbrmarkers\tikzset{mkl/.style={align=left}}\else\tikzset{mkl/.style={}}\fi
  % ---- column headers, with arrows from step to step --------------------------------
  \node[hdr,fill=psblue!75!black,minimum width=2.6cm] at (0,0)
    {1 configure\\[-1pt]{\mdseries cmake reads the CPU}};
  \node[hdr,fill=psblue!75!black,minimum width=4.1cm] at (2.85,0)
    {2 generate code\\[-1pt]{\mdseries make runs Python}};
  \node[hdr,fill=psblue!75!black,minimum width=3.4cm] at (7.2,0)
    {3 compile\\[-1pt]{\mdseries make runs GCC or Clang}};
  \node[hdr,fill=psteal!80!black,minimum width=3.55cm] at (10.85,0)
    {4 run\\[-1pt]{\mdseries no Python}};
  \foreach \x in {2.62,6.97,10.62}
    \fill[hdrarr] (\x,-0.22) -- ++(0.2,-0.18) -- ++(-0.2,-0.18) -- cycle;

  % ---- 1 configure: two settings for the later steps --------------------------------
  \node[b,text width=2.4cm] (a3) at (1.3,-1.4) {\textbf{target} for step 2:\\RVV with this CPU's \textit{V}};
  \node[b,text width=2.4cm,font=\tiny] (a2) at (1.3,-2.55)
    {{\scriptsize\textbf{flags} for step 3:}\\\code{-march=rv64gcv\_zvl}\textit{V}\code{b}};
  \node[anchor=north west,font=\tiny,align=left,inner sep=0pt] at (0.1,-3.55)
    {\tikz\fill[psviolet!20,draw=psviolet!80!black] (0,0) rectangle (2.4mm,1.6mm);\ generated code\\[3pt]
     \tikz\draw[psgray!80] (0,0) rectangle (2.4mm,1.6mm);\ written by hand};

  % ---- 2 generate: the Python tool chain ---------------------------------------------
  \node[b,text width=3.85cm] (g0) at (4.9,-1.27) {\textbf{generator script} of the program};
  \node[b,text width=3.85cm] (g1) at (4.9,-2.17) {\textbf{lbmpy}: formulas for one cell};
  \node[b,text width=3.85cm] (g2) at (4.9,-3.2) {\textbf{pystencils}: loop over all cells,\\C++ with RVV};
  \node[b,text width=3.85cm] (g3) at (4.9,-4.35) {\textbf{sweepgen}: writes the C++ files,\\a class around the loop};
  \foreach \u/\v in {g0/g1,g1/g2,g2/g3} \draw[arr] (\u) -- (\v);

  % ---- 3 compile: hand-written and generated C++ together ----------------------------
  \node[b,mkl,text width=3.15cm] (i1) at (8.9,-1.27) {\textbf{waLBerla core} library};
  \node[b,mkl,text width=3.15cm] (i2) at (8.9,-1.92) {the program's \code{main()}};
  \node[gen,text width=3.15cm] (i0) at (8.9,-2.57) {generated C++ files};
  \begin{scope}[on background layer]
    \node[fill=psgray!12,rounded corners=2pt,inner sep=2.5pt,fit=(i1)(i2)(i0)] (src) {};
  \end{scope}
  \node[b,text width=3.15cm,fill=psgray!8] (cc) at (8.9,-3.6) {\textbf{GCC or Clang}\\with the flags of step 1};
  \node[b,text width=3.15cm] (exe) at (8.9,-4.47) {the executable};
  \draw[arr] (src) -- (cc);
  \draw[arr] (cc) -- (exe);

  % ---- 4 run: set up, then the time loop (kernel, halo exchange, wall pass; notes5 F16) ----
  \node[b,text width=2.8cm] (r0) at (12.45,-1.2) {set up blocks and fields};
  \node[gen,text width=2.8cm] (r1) at (12.45,-2.05) {\textbf{generated kernel}\\on each block};
  \node[b,text width=2.8cm] (r2) at (12.45,-2.88) {\textbf{halo exchange}};
  \node[gen,text width=2.8cm] (rw) at (12.45,-3.6) {\textbf{wall pass}, plain C++};
  \node[b,text width=2.8cm] (r3) at (12.45,-4.45) {print MLUP/s and checksum};
  \draw[arr] (r0) -- (r1);
  \draw[arr] (r1) -- (r2);
  \draw[arr] (r2) -- (rw);
  \draw[arr] (rw.east) -- ++(0.12,0) |- (r1.east);
  \node[lbl,rotate=90] at (14.27,-2.82) {200 steps};
  \draw[arr] (rw) -- (r3);

  % ---- time bands ----------------------------------------------------------------------
  \fill[psblue!8] (0,-4.95) rectangle (10.6,-5.37);
  \node[font=\scriptsize,text=psblue!75!black] at (5.3,-5.16)
    {\textbf{build time}: once, on the board (all builds together take hours)};
  \fill[psteal!10] (10.85,-4.95) rectangle (14.4,-5.37);
  \node[font=\scriptsize,text=psteal!75!black] at (12.62,-5.16) {\textbf{run time}: seconds};

  % ---- markers: what we changed ------------------------------------------------------
  \ifbrmarkers
    \tikzset{mk/.append style={minimum size=3.2mm}}
    %   \brmark: on top of the box's upper border; \brmarkin: inside the box, at its right end
    \def\brmark#1#2#3#4{% node, slot (0 = rightmost), colour, number
      \node[mk,fill=#3] at ([xshift={-2.2mm-#2*3.6mm},yshift=1.3mm]#1.north east) {#4};}
    \def\brmarkin#1#2#3#4{%
      \node[mk,fill=#3] at ([xshift={-2.2mm-#2*3.6mm}]#1.east) {#4};}
    \brmark{a2}{0}{psblue}{1}   \brmark{a2}{1}{psteal}{8}
    \brmark{g0}{0}{psblue}{3}
    \brmark{g2}{0}{psblue}{2}   \brmark{g2}{1}{psteal}{5}   \brmark{g2}{2}{psteal}{12}
    \brmark{g3}{0}{psteal}{6}   \brmark{g3}{1}{psteal}{10}
    \brmarkin{i1}{0}{psteal}{11} \brmarkin{i1}{1}{psteal}{7}
    \brmarkin{i2}{0}{psteal}{9}    \brmarkin{i2}{1}{psblue}{3}
    \brmark{cc}{0}{psblue}{4}
  \fi
\end{tikzpicture}
